El 30\,\% dels clients d'una botiga en línia són socis del seu club ($S$). El 70\,\% dels socis fan més d'una
compra al mes ($M$), mentre que, entre els que no en són socis, només ho fa el 20\,\%. Es tria un client a
l'atzar.

\begin{apartats}

\apartat[1,25]{0,75}
Quina és la probabilitat que faci més d'una compra al mes?

\begin{solucio}
$P(M)=P(S)\,P(M\mid S)+P\left(\overline{S}\right)P\left(M\mid\overline{S}\right)=0{,}3\cdot0{,}7+0{,}7\cdot0{,}2
=0{,}21+0{,}14=0{,}35$.
\end{solucio}

\begin{tria}{bayes-tasca}
\itemtria{original}{1}{1,25}
Si fa més d'una compra al mes, quina és la probabilitat que sigui soci? I si no en fa més d'una?

\begin{solucio}
$P(S\mid M)=\dfrac{0{,}21}{0{,}35}=0{,}6$.\\
$P\left(S\mid\overline{M}\right)=\dfrac{0{,}3\cdot0{,}3}{1-0{,}35}=\dfrac{0{,}09}{0{,}65}\approx0{,}138$.
\end{solucio}

\itemtria{arbre-invers}{1}{1,25}
Construeix l'arbre en l'ordre invers: primer, si el client fa més d'una compra al mes, i després, si és
soci. Calcula totes les probabilitats de les branques.

\begin{solucio}
Primer nivell: $P(M)=0{,}35$ i $P\left(\overline{M}\right)=0{,}65$.\\
Segon nivell, amb el teorema de Bayes: $P(S\mid M)=\dfrac{0{,}21}{0{,}35}=0{,}6$ i
$P\left(\overline{S}\mid M\right)=0{,}4$; $P\left(S\mid\overline{M}\right)=\dfrac{0{,}09}{0{,}65}\approx0{,}138$
i $P\left(\overline{S}\mid\overline{M}\right)=\dfrac{0{,}56}{0{,}65}\approx0{,}862$.\\
Comprovació: els camins donen les mateixes interseccions que l'arbre original; per exemple,
$0{,}35\cdot0{,}6=0{,}21=P(S\cap M)$.
\end{solucio}
\end{tria}

\begin{nomesllarg}
\apartat{0,75}
Són independents «ser soci» i «fer més d'una compra al mes»? Justifica-ho.

\begin{solucio}
$P(M\mid S)=0{,}7$, però $P(M)=0{,}35$. Com que saber que és soci canvia la probabilitat, \textbf{no} són
independents. (També: $P(S\cap M)=0{,}21\neq P(S)\cdot P(M)=0{,}105$.)
\end{solucio}
\end{nomesllarg}

\end{apartats}
